\section{Consequences}\label{sec:consequences}

The reduction and fractional construction now give the unconditional
configuration gap. We then use a satisfiability gap to obtain the
algorithmic hardness result.

\subsection{The unconditional configuration gap}

\begin{proof}[Proof of Theorem~\ref{thm:unbounded-gap}]
Fix an integer $c\ge0$, set $\rho=1/8$, and apply the actual construction
of Theorem~\ref{thm:reduction} to the four-vertex clique with $k=2$.
All parameters, including $d$, the trees, and
$R=100(P+4K+1)$, are chosen as prescribed for these constants.
The minimum vertex cover of this graph has size three, whereas
\begin{equation}\label{eq:clique-cover-contradiction}
  k+\rho n=2+\frac48=\frac52<3.
\end{equation}
Proposition~\ref{prop:soundness} therefore rules out every packing into
at most $B+c$ bins.  Independently,
Proposition~\ref{prop:fractional} gives configuration-LP value exactly
$B$, for both individual-copy and numerical-type configurations.
This proves Equation~\eqref{eq:main-gap} for every $c$.
\end{proof}

The structural proof uses only the four-vertex clique, the finite
construction, the fractional witness, and the soundness bound. The PCP
Theorem enters only the algorithmic hardness argument below.

\subsection{A vertex-cover gap with perfect completeness}

We use the following standard consequence of the PCP Theorem.
For some absolute constant $\eta>0$, it is NP-hard to distinguish
satisfiable formulas having exactly three literal slots per clause from
formulas in which no assignment satisfies more than a $1-\eta$ fraction
of the clauses.  The clause list may be required to be nonempty.
H\aa stad's Theorem~6.5~\cite{Hastad2001} gives this statement with
$\eta=1/16$ by taking $\varepsilon=1/16$ in its
$7/8+\varepsilon$ soundness bound.  In particular the YES case is fully
satisfiable, as needed for the completeness direction of our reduction.

We use a literal-slot version of the complement of the
satisfiability-to-clique graph described by Karp~\cite[p.~97]{Karp1972}.
Given such a formula with $M\ge1$ clauses, create three vertices for its
literal slots in each clause and join those vertices into a triangle.
Also join any two slots occupied by complementary literals, retaining
only one edge when an edge is specified more than once.  The result is
a simple graph with
\begin{equation}\label{eq:clause-graph-parameters}
  n=3M,\qquad k=2M=2n/3.
\end{equation}
Every vertex lies in a triangle, so the graph has no isolated vertices
and at least one edge.  The construction is polynomial in the explicit
formula encoding.

An independent set selects at most one slot from any clause and never
selects complementary literals.  Setting its selected literals true and
extending the assignment arbitrarily satisfies at least as many clauses
as the cardinality of the independent set.  Conversely, from any assignment
choose one true slot in each satisfied clause.  These slots form an
independent set.  Thus, writing $\operatorname{MAXSAT}(F)$ for the largest
number of simultaneously satisfiable clauses, $\alpha(G)$ for the
maximum independent-set size, and $\tau(G)$ for the minimum vertex-cover
size,
\begin{equation}\label{eq:clause-graph-identity}
  \alpha(G)=\operatorname{MAXSAT}(F),
  \qquad \tau(G)=3M-\operatorname{MAXSAT}(F).
\end{equation}
This argument also permits repeated literals and clauses.
The minimum cover is $k$ in the satisfiable case and at least
\begin{equation}\label{eq:vertex-cover-gap}
  (2+\eta)M=k+\frac{\eta}{3}n
\end{equation}
in the other case.

\subsection{Additive hardness and the complexity equivalence}

\begin{proof}[Proof of Theorem~\ref{thm:additive-hardness}]
Fix $c\ge0$ and fix $0<\rho<\eta/3$, for example $\rho=1/96$ with
the preceding choice of $\eta$.  Apply Theorem~\ref{thm:reduction} to
the clause graph and target in Equation~\eqref{eq:clause-graph-parameters}.
The resulting rational list and integer $B$ have polynomial encoding
length and are computable in polynomial time for these fixed $c,\rho$.

In the satisfiable case the graph has a cover of size $k$, and
Proposition~\ref{prop:completeness} gives a packing into $B$ bins.
In the other case, a packing into at most $B+c$ bins would, by
Proposition~\ref{prop:soundness}, give a cover of size at most
$k+\rho n$.  This contradicts Equation~\eqref{eq:vertex-cover-gap}.
All constructed sizes exceed $1/6$.  The promised distinction is
therefore NP-hard.
\end{proof}

\begin{proof}[Proof of Corollary~\ref{cor:complexity-equiv}]
Suppose such an algorithm exists, and fix a nonnegative integer $c$
at least as large as its additive constant $C$.
On a YES instance of Theorem~\ref{thm:additive-hardness}, the algorithm
must return at most $\OPT(I)+C\le B+c$ bins.
On a NO instance, no packing using that many bins exists.
Counting the bins in the returned packing distinguishes the two cases
in deterministic polynomial time.  Hence $\classP=\classNP$.

Conversely, assume $\classP=\classNP$.
For an explicit rational list of $N$ items and a bound $b\le N$,
bin feasibility belongs to $\classNP$: a certificate assigns each item
a bin index, and every capacity constraint can be checked with exact
rational arithmetic in polynomial bit time.
The same statement holds when some item assignments have already been
fixed.  Under the assumption these decision problems are solvable in
polynomial time.  Binary search over $0\le b\le N$ finds the optimum.
Then assign items successively, trying at most $b$ bin indices at each
step and retaining one for which a completion remains feasible.
At most $N^2$ further tests construct an optimal packing.  The empty
list is handled directly.  This gives the desired algorithm with $C=0$.
\end{proof}
